%! Logarithmic Expressions — Unit 5, Lesson 5.1 — Group Activity
\documentclass[10pt]{article}
\usepackage{precalculus-article}
\usepackage{precalculus-boxes}
\newcommand{\TallMath}[1]{$\displaystyle #1\rule[-1.4em]{0pt}{3.2em}$}

\begin{document}

\pageheader{Unit 5, Lesson 5.1}{Group Activity}
\namepartnerperiod

\vspace{0.08in}

\begin{scenariobox}[Logarithms in the Real World]{plumlight}
\small
Scientists use logarithms to measure quantities that span huge ranges. The \textbf{decibel scale}
measures sound intensity: a whisper is about 30 dB and a jet engine about 140 dB, yet the actual
intensities differ by a factor of $10^{11}$. The \textbf{pH scale} measures acidity: each unit
of pH corresponds to a 10$\times$ change in hydrogen-ion concentration. Work through the tier that
best matches your current level.
\end{scenariobox}

\vspace{0.08in}

\begin{tcolorbox}[colback=white, colframe=black!40,
    title=\textbf{Tier R --- Remediate},
    fonttitle=\bfseries, arc=2mm, left=3mm, right=3mm, top=2mm, bottom=2mm]

\textbf{Part A: Conversion drill.} Convert each equation to the other form.

\renewcommand{\arraystretch}{1.9}
\begin{tabularx}{\linewidth}{@{} X X @{}}
\textbf{Exponential $\to$ Logarithmic} & \textbf{Logarithmic $\to$ Exponential} \\
\midrule
$3^4 = 81$ \quad $\Rightarrow$ \blank{3.5cm}  & $\log_2 32 = 5$ \quad $\Rightarrow$ \blank{3.5cm} \\[4pt]
$5^3 = 125$ \quad $\Rightarrow$ \blank{3.5cm} & $\log_3 9 = 2$ \quad $\Rightarrow$ \blank{3.5cm} \\[4pt]
$10^2 = 100$ \quad $\Rightarrow$ \blank{3.5cm} & $\log_6 1 = 0$ \quad $\Rightarrow$ \blank{3.5cm} \\[4pt]
$4^{1/2} = 2$ \quad $\Rightarrow$ \blank{3.5cm} & $\log_{10} 0.1 = -1$ \quad $\Rightarrow$ \blank{3.5cm} \\
\end{tabularx}

\vspace{0.10in}
\textbf{Part B: Evaluate each logarithm exactly.}

\vspace{0.06in}
$\log_2 16 = $ \blank{2cm} \qquad
$\log_3 1 = $ \blank{2cm} \qquad
$\log_{10} 10{,}000 = $ \blank{2cm} \qquad
$\log_5 5 = $ \blank{2cm}

\end{tcolorbox}

\vspace{0.08in}

\begin{tcolorbox}[colback=white, colframe=black!40,
    title=\textbf{Tier A --- Apply},
    fonttitle=\bfseries, arc=2mm, left=3mm, right=3mm, top=2mm, bottom=2mm]

\begin{enumerate}[itemsep=10pt]

  \item Evaluate each logarithm exactly. Show the exponential form you used.
        \begin{enumerate}[label=(\alph*), itemsep=4pt]
          \item $\log_6 36 = $ \blank{2.5cm} \hspace{1em} because $6^{\blank{1cm}} = 36$
          \item $\log_{10} 0.001 = $ \blank{2.5cm} \hspace{1em} because $10^{\blank{1cm}} = 0.001$
          \item $\log_9 3 = $ \blank{2.5cm} \hspace{1em} because $9^{\blank{1cm}} = 3$
        \end{enumerate}

  \item Use a calculator to estimate to 3 decimal places.
        \begin{enumerate}[label=(\alph*), itemsep=4pt]
          \item $\log_2 7 \approx $ \blank{3cm}
          \item $\ln 5 \approx $ \blank{3cm}
        \end{enumerate}

  \item Without a calculator, determine between which two consecutive integers $\log_3 20$ falls.
        Explain your reasoning.

        \vspace{0.04in}
        $3^{\blank{1cm}} = $ \blank{1.5cm} $< 20 <$ \blank{1.5cm} $= 3^{\blank{1cm}}$,
        so $\log_3 20$ is between \blank{1cm} and \blank{1cm}.

  \item The pH of a solution is $\text{pH} = -\log_{10}[\text{H}^+]$, where $[\text{H}^+]$
        is the hydrogen-ion concentration in mol/L.
        \begin{enumerate}[label=(\alph*), itemsep=4pt]
          \item If $[\text{H}^+] = 10^{-4}$, find the pH. \blank{3cm}
          \item If $[\text{H}^+] = 10^{-7}$, find the pH. \blank{3cm}
          \item Which solution is more acidic (lower pH)? \blank{3cm}
        \end{enumerate}

\end{enumerate}
\end{tcolorbox}

\vspace{0.08in}

\begin{tcolorbox}[colback=white, colframe=black!40,
    title=\textbf{Tier E --- Extend},
    fonttitle=\bfseries, arc=2mm, left=3mm, right=3mm, top=2mm, bottom=2mm]

\begin{enumerate}[itemsep=10pt]

  \item \textbf{Prove two inverse properties} using only the definition
        $\log_b c = a \Leftrightarrow b^a = c$.

        \begin{enumerate}[label=(\alph*), itemsep=8pt]
          \item Show that $\log_b b^n = n$ for any real $n$.

                \writelines{3}

          \item Show that $b^{\log_b n} = n$ for any $n > 0$.

                \writelines{3}
        \end{enumerate}

  \item Solve for $x$: $\log_4 x = \dfrac{3}{2}$.

        Convert to exponential form, then simplify:

        \vspace{0.06in}
        $4^{3/2} = $ \blank{3cm}, so $x = $ \blank{2.5cm}

  \item Explain in 2--3 sentences: what is the connection between a \textbf{fractional exponent}
        (such as $4^{3/2}$) and a \textbf{fractional logarithm} (such as $\log_4 8 = 3/2$)?

        \writelines{3}

\end{enumerate}
\end{tcolorbox}

\end{document}
